\nonstopmode{}
\documentclass[letterpaper]{book}
\usepackage[times,inconsolata,hyper]{Rd}
\usepackage{makeidx}
\makeatletter\@ifl@t@r\fmtversion{2018/04/01}{}{\usepackage[utf8]{inputenc}}\makeatother
% \usepackage{graphicx} % @USE GRAPHICX@
\makeindex{}
\begin{document}
\chapter*{}
\begin{center}
{\textbf{\huge Package `stimgate'}}
\par\bigskip{\large \today}
\end{center}
\ifthenelse{\boolean{Rd@use@hyper}}{\hypersetup{pdftitle = {stimgate: Identify responding cells as outliers}}}{}
\ifthenelse{\boolean{Rd@use@hyper}}{\hypersetup{pdfauthor = {Miguel Rodo}}}{}
\begin{description}
\raggedright{}
\item[Type]\AsIs{Package}
\item[Title]\AsIs{Identify responding cells as outliers}
\item[Version]\AsIs{0.99.2-4}
\item[Description]\AsIs{Apply outlier gating function to identify cells that have possibly
responded to stimulation, by comparing the unstimulated and stimulated tubes
from the same sample.}
\item[URL]\AsIs{}\url{https://github.com/SATVILab/stimgate}\AsIs{}
\item[BugReports]\AsIs{}\url{https://github.com/SATVILab/stimgate/issues}\AsIs{}
\item[License]\AsIs{GPL (>= 2)}
\item[Encoding]\AsIs{UTF-8}
\item[Depends]\AsIs{R (>= 4.4.0)}
\item[Roxygen]\AsIs{list(markdown = TRUE)}
\item[RoxygenNote]\AsIs{7.3.3}
\item[biocViews]\AsIs{FlowCytometry, SingleCell, CellBasedAssays, ImmunoOncology}
\item[VignetteEngine]\AsIs{knitr::rmarkdown}
\item[Imports]\AsIs{cowplot, dplyr, ggplot2, purrr, rlang, stringr, tibble, tidyr,
scam, mgcv, gtools, flowCore, flowWorkspace, cluster}
\item[Suggests]\AsIs{testthat (>= 3.0.0), HDCytoData, here, BiocManager, knitr,
rmarkdown, hexbin}
\item[NeedsCompilation]\AsIs{no}
\item[Author]\AsIs{Miguel Rodo [aut, cre] (ORCID: <}\url{https://orcid.org/0000-0002-2036-4878}\AsIs{>)}
\item[Maintainer]\AsIs{Miguel Rodo }\email{miguel.rodo@uct.ac.za}\AsIs{}
\end{description}
\Rdcontents{Contents}
\HeaderA{.assert\_string}{Assert that object is a non-empty string}{.assert.Rul.string}
\keyword{internal}{.assert\_string}
%
\begin{Description}
Assert that object is a non-empty string
\end{Description}
%
\begin{Usage}
\begin{verbatim}
.assert_string(x)
\end{verbatim}
\end{Usage}
%
\begin{Arguments}
\begin{ldescription}
\item[\code{x}] object Object to check
\end{ldescription}
\end{Arguments}
%
\begin{Value}
invisible NULL if check passes, otherwise throws an error
\end{Value}
\HeaderA{.assert\_string\_vector}{Assert that object is a non-empty character vector}{.assert.Rul.string.Rul.vector}
\keyword{internal}{.assert\_string\_vector}
%
\begin{Description}
Assert that object is a non-empty character vector
\end{Description}
%
\begin{Usage}
\begin{verbatim}
.assert_string_vector(x)
\end{verbatim}
\end{Usage}
%
\begin{Arguments}
\begin{ldescription}
\item[\code{x}] object Object to check
\end{ldescription}
\end{Arguments}
%
\begin{Value}
invisible NULL if check passes, otherwise throws an error
\end{Value}
\HeaderA{.cytokine\_cutpoint}{Constructs a cutpoint for a flowFrame by using a derivative of the kernel density estimate}{.cytokine.Rul.cutpoint}
%
\begin{Description}
We determine a gating cutpoint using either the first or second derivative of
the kernel density estimate (KDE) of the \code{x}.
\end{Description}
%
\begin{Usage}
\begin{verbatim}
.cytokine_cutpoint(
  x,
  num_peaks = 1,
  ref_peak = 1,
  method = c("first_deriv", "second_deriv"),
  tol = 0.01,
  adjust = 1,
  side = "right",
  strict = TRUE,
  plot = FALSE,
  auto_tol = FALSE,
  ...
)
\end{verbatim}
\end{Usage}
%
\begin{Arguments}
\begin{ldescription}
\item[\code{x}] a \code{numeric} vector used as input data

\item[\code{num\_peaks}] the number of peaks expected to see. This effectively removes
any peaks that are artifacts of smoothing

\item[\code{ref\_peak}] After \code{num\_peaks} are found, this argument provides the
index of the reference population from which a gate will be obtained. By
default, the peak farthest to the left is used.

\item[\code{method}] the method used to select the cutpoint. See details.

\item[\code{tol}] the tolerance value

\item[\code{adjust}] the scaling adjustment applied to the bandwidth used in the
first derivative of the kernel density estimate

\item[\code{side}] character specifying the side of the gate, either
\code{'right'} (default) or \code{'left'}

\item[\code{strict}] \code{logical} when the actual number of peaks detected is less than \code{ref\_peak}.
an error is reported by default. But if \code{strict} is set to FALSE, then the reference peak will be reset to the peak of the far right.

\item[\code{plot}] logical specifying whether to plot the peaks found

\item[\code{auto\_tol}] when TRUE, it tries to set the tolerance automatically.

\item[\code{...}] additional arguments passed to \code{.deriv\_density}
\end{ldescription}
\end{Arguments}
%
\begin{Details}
By default, we compute the first derivative of the kernel density estimate.
Next, we determine the lowest valley from the derivative, which corresponds to the
density's mode for cytokines. We then contruct a gating cutpoint as the value
less than the tolerance value \code{tol} in magnitude and is also greater
than the lowest valley.

Alternatively, if the \code{method} is selected as \code{second\_deriv}, we
select a cutpoint from the second derivative of the KDE. Specifically, we
choose the cutpoint as the largest peak of the second derivative of the KDE
density which is greater than the reference peak.
\end{Details}
%
\begin{Value}
the cutpoint along the x-axis
\end{Value}
\HeaderA{.debug\_msg}{Print debug message conditionally}{.debug.Rul.msg}
\keyword{internal}{.debug\_msg}
%
\begin{Description}
Print debug message conditionally
\end{Description}
%
\begin{Usage}
\begin{verbatim}
.debug_msg(.debug, msg, val = NULL)
\end{verbatim}
\end{Usage}
%
\begin{Arguments}
\begin{ldescription}
\item[\code{.debug}] logical Whether to print debug messages

\item[\code{msg}] character Message to print

\item[\code{val}] object Optional value to append to message. Default is NULL.
\end{ldescription}
\end{Arguments}
%
\begin{Value}
logical invisibly TRUE if message was printed, FALSE otherwise
\end{Value}
\HeaderA{.is\_string}{Check if object is a non-empty string}{.is.Rul.string}
\keyword{internal}{.is\_string}
%
\begin{Description}
Check if object is a non-empty string
\end{Description}
%
\begin{Usage}
\begin{verbatim}
.is_string(x)
\end{verbatim}
\end{Usage}
%
\begin{Arguments}
\begin{ldescription}
\item[\code{x}] object Object to check
\end{ldescription}
\end{Arguments}
%
\begin{Value}
logical TRUE if x is a character vector of length 1 with non-zero characters
\end{Value}
\HeaderA{.is\_string\_vector}{Check if object is a non-empty character vector}{.is.Rul.string.Rul.vector}
\keyword{internal}{.is\_string\_vector}
%
\begin{Description}
Check if object is a non-empty character vector
\end{Description}
%
\begin{Usage}
\begin{verbatim}
.is_string_vector(x)
\end{verbatim}
\end{Usage}
%
\begin{Arguments}
\begin{ldescription}
\item[\code{x}] object Object to check
\end{ldescription}
\end{Arguments}
%
\begin{Value}
logical TRUE if x is a character vector of length >= 1 with all non-zero character elements
\end{Value}
\HeaderA{chnl\_lab}{Get markers and channels}{chnl.Rul.lab}
\aliasA{chnl\_to\_marker,}{chnl\_lab}{chnl.Rul.to.Rul.marker,}
\aliasA{get\_chnl}{chnl\_lab}{get.Rul.chnl}
\aliasA{get\_marker,}{chnl\_lab}{get.Rul.marker,}
\aliasA{marker\_lab,}{chnl\_lab}{marker.Rul.lab,}
\aliasA{marker\_to\_chnl,}{chnl\_lab}{marker.Rul.to.Rul.chnl,}
%
\begin{Description}
From a cytometry object (e.g. flowFrame or flowSet),
either get a character vector of markers
or channels (get\_chnl and get\_marker),
or get a named vector that converts
between channel names and marker names (e.g. chnl\_to\_marker).
\end{Description}
%
\begin{Usage}
\begin{verbatim}
chnl_lab(data)
\end{verbatim}
\end{Usage}
%
\begin{Arguments}
\begin{ldescription}
\item[\code{data}] object of class flowFrame, flowSet. Channel and corresponding
marker names are drawn from here.
\end{ldescription}
\end{Arguments}
%
\begin{Details}
Note that chnl\_lab is equivalent to chnl\_to\_marker,
and marker\_lab is equivalent to marker\_to\_chnl.
\end{Details}
%
\begin{Value}
A named character vector.
\end{Value}
%
\begin{Examples}
\begin{ExampleCode}

# Create example flowFrame-like data structure
data(GvHD, package = "flowCore") 
fs <- GvHD[1:2]

# Get channel to marker mapping
chnl_lab(fs[[1]])


\end{ExampleCode}
\end{Examples}
\HeaderA{get\_example\_data}{Get example GatingSet}{get.Rul.example.Rul.data}
%
\begin{Description}
Create and save a complete example GatingSet for testing and examples.
This function internally creates a flowSet, samples channels, and saves the GatingSet.
\end{Description}
%
\begin{Usage}
\begin{verbatim}
get_example_data(dir_cache = NULL)
\end{verbatim}
\end{Usage}
%
\begin{Arguments}
\begin{ldescription}
\item[\code{dir\_cache}] Directory to save the GatingSet. If NULL, uses a temporary directory.
\end{ldescription}
\end{Arguments}
%
\begin{Value}
A list containing the path to the saved GatingSet, batch\_list, and marker names
\end{Value}
\HeaderA{get\_stats}{Get gating statistics}{get.Rul.stats}
%
\begin{Description}
Get gating statistics
\end{Description}
%
\begin{Usage}
\begin{verbatim}
get_stats(path_project)
\end{verbatim}
\end{Usage}
%
\begin{Arguments}
\begin{ldescription}
\item[\code{path\_project}] character. Path to the project directory.
\end{ldescription}
\end{Arguments}
%
\begin{Value}
A data frame with gating statistics.
\end{Value}
%
\begin{Examples}
\begin{ExampleCode}
{
# Get example dataset
example_data <- get_example_data()
gs <- flowWorkspace::load_gs(example_data$path_gs)

# Run the stimgate pipeline
path_project <- stimgate_gate(
  path_project = file.path(tempdir(), "stimgate_example"),
  .data = gs,
  batch_list = example_data$batch_list,
  marker = example_data$marker,
  pop_gate = "root"
)

# Get statistics for the identified gates
stats <- get_stats(path_project)
}
\end{ExampleCode}
\end{Examples}
\HeaderA{stimgate\_data\_get\_ex}{Read saved expression data from project}{stimgate.Rul.data.Rul.get.Rul.ex}
%
\begin{Description}
Read channel expression vectors saved under a project's sample\_data
directory and return them as a tibble with sample metadata columns.
\end{Description}
%
\begin{Usage}
\begin{verbatim}
stimgate_data_get_ex(
  path_project,
  pop = NULL,
  ind = NULL,
  chnl = NULL,
  bias = FALSE,
  exc_min = FALSE,
  combn_exc = NULL,
  cyt_pos = FALSE,
  gate_type_cyt_pos = "cyt",
  gate_type_single_pos = "single",
  mult = FALSE,
  gate_uns_method = "min"
)
\end{verbatim}
\end{Usage}
%
\begin{Arguments}
\begin{ldescription}
\item[\code{path\_project}] character Path to project.

\item[\code{pop}] character or NULL Population name(s). Default is detected from project sample\_data.

\item[\code{ind}] character or NULL Index/indices of samples. Default is detected from project sample\_data.

\item[\code{chnl}] character or NULL Channel name(s) to return. Default is detected from project sample\_data.

\item[\code{bias}] logical. Whether to add bias to unstimulated sample used in the gating. Default is \code{FALSE}.

\item[\code{exc\_min}] logical. Whether to exclude cells with the minimum expression for any channels.

\item[\code{combn\_exc}] list. Combinations of channels to exclude. Default is NULL.

\item[\code{cyt\_pos}] logical. Whether to return only cytokine-positive cells. Default is FALSE.

\item[\code{gate\_type\_cyt\_pos}] character. Gate type to use for cytokine-positive cells. Default is "cyt".

\item[\code{gate\_type\_single\_pos}] character. Gate type to use for single-positive cells. Default is "single".

\item[\code{mult}] logical. Whether to return only multi-functional cells (positive for multiple markers). Default is FALSE.

\item[\code{gate\_uns\_method}] character. Method for gating unstimulated cells. Default is "min".
\end{ldescription}
\end{Arguments}
%
\begin{Value}
A tibble with columns \code{pop}, \code{ind} and one column per requested channel. Rows correspond to cells.
\end{Value}
%
\begin{Examples}
\begin{ExampleCode}
## Not run: 
tmp <- tempdir()
dir.create(file.path(tmp, "sample_data", "POP1", "ind_1"), recursive = TRUE)
saveRDS(c(1, 2, 3), file.path(tmp, "sample_data", "POP1", "ind_1", "chnl_BC1.rds"))
saveRDS(c(4, 5, 6), file.path(tmp, "sample_data", "POP1", "ind_1", "chnl_BC2.rds"))
stimgate_data_get_ex(tmp)
stimgate_data_get_ex(tmp, chnl = "BC1")

## End(Not run)
\end{ExampleCode}
\end{Examples}
\HeaderA{stimgate\_fcs\_write}{Write FCS files of marker-positive FCS files}{stimgate.Rul.fcs.Rul.write}
%
\begin{Description}
Uses the gates to write FCS files of marker-positive FCS files.
\end{Description}
%
\begin{Usage}
\begin{verbatim}
stimgate_fcs_write(
  path_project,
  .data,
  pop = NULL,
  ind_batch_list,
  path_dir_save,
  chnl = NULL,
  gate_tbl = NULL,
  trans_fn = NULL,
  trans_chnl = NULL,
  combn_exc = NULL,
  gate_type_cyt_pos = "cyt",
  gate_type_single_pos = "single",
  mult = FALSE,
  gate_uns_method = "min"
)
\end{verbatim}
\end{Usage}
%
\begin{Arguments}
\begin{ldescription}
\item[\code{path\_project}] character. Path to project directory.

\item[\code{.data}] GatingSet. GatingSet object containing the flow cytometry data.

\item[\code{pop}] character. Population that was gated on.

\item[\code{ind\_batch\_list}] list. List of indices grouped by batch.

\item[\code{path\_dir\_save}] character. Directory path to save the FCS files to.

\item[\code{chnl}] character vector. Specific channels to gate on.

\item[\code{gate\_tbl}] data.frame. Pre-computed gate table, if available.

\item[\code{trans\_fn}] function. Transformation function to apply.

\item[\code{trans\_chnl}] character vector. Columns to transform.

\item[\code{combn\_exc}] list. Combinations of channels to exclude.

\item[\code{gate\_type\_cyt\_pos}] character. Gate type to use for cytokine-positive cells.

\item[\code{gate\_type\_single\_pos}] character. Gate type to use for single-positive cells.

\item[\code{mult}] logical. Whether cells must be multi-positive.

\item[\code{gate\_uns\_method}] character. Method to calculate unstimulated thresholds.
\end{ldescription}
\end{Arguments}
%
\begin{Details}
This function processes flow cytometry data to identify and export cytokine-positive
cells to FCS files. It requires that gates have been pre-computed using
\code{\LinkA{stimgate\_gate}{stimgate.Rul.gate}} or that a complete gate table is provided.

The function will create the output directory and write FCS files for samples
that contain cytokine-positive cells. If no positive cells are found in a sample,
no FCS file will be written for that sample.
\end{Details}
%
\begin{Examples}
\begin{ExampleCode}
## Not run: 
# Complete workflow example
library(stimgate)

# Load your GatingSet (gs) and define batch structure
# batch_list <- list(batch1 = c(1, 2, 3), batch2 = c(4, 5, 6))
# where the last element in each batch is the unstimulated sample

# First, run gating to create gates
path_project <- tempfile("stimgate_project")
# stimgate_gate(
#   .data = gs,
#   path_project = path_project,
#   pop_gate = "root",
#   batch_list = batch_list,
#   marker = c("IL2", "IFNg")  # your cytokine markers
# )

# Then write FCS files of cytokine-positive cells
path_output <- tempfile("fcs_output")
# stimgate_fcs_write(
#   path_project = path_project,
#   .data = gs,
#   ind_batch_list = batch_list,
#   path_dir_save = path_output,
#   chnl = c("IL2", "IFNg")
# )

# Alternative: provide your own gate table
# gate_tbl <- data.frame(
#   chnl = c("IL2", "IFNg"),
#   marker = c("IL2", "IFNg"),
#   batch = c(1, 1),
#   ind = c(1, 1),
#   gate = c(0.5, 0.3),
#   gate_name = c("gate", "gate")
# )
# stimgate_fcs_write(
#   path_project = path_project,
#   .data = gs,
#   ind_batch_list = batch_list,
#   path_dir_save = path_output,
#   chnl = c("IL2", "IFNg"),
#   gate_tbl = gate_tbl
# )

## End(Not run)
\end{ExampleCode}
\end{Examples}
\HeaderA{stimgate\_gate}{Identify cytokine-positive cells through automated gating}{stimgate.Rul.gate}
%
\begin{Description}
Main function for identifying cytokine-positive cells using outlier-based gating
to compare stimulated versus unstimulated samples. This function implements a
comprehensive workflow that identifies cells responding to stimulation by
detecting outliers in cytokine expression distributions. The process includes
density estimation, threshold identification, clustering-based gate refinement,
and generation of comprehensive statistics and visualizations.

The function operates by comparing cytokine expression in stimulated samples
against corresponding unstimulated controls from the same donor/batch to identify
cells that have likely responded to stimulation. It accounts for batch effects,
background cytokine production, and technical variability.
\end{Description}
%
\begin{Usage}
\begin{verbatim}
stimgate_gate(
  path_project,
  .data,
  batch_list,
  chnl = NULL,
  marker = NULL,
  pop_gate = "root",
  bias_uns = NULL,
  bias_uns_factor = 1,
  exc_min = TRUE,
  cp_min = NULL,
  bw_min = NULL,
  min_cell = 100,
  max_pos_prob_x = Inf,
  gate_quant = c(0.25, 0.75),
  tol_clust = 1e-07,
  gate_combn = "min",
  marker_settings = NULL,
  chnl_settings = NULL,
  calc_cyt_pos_gates = TRUE,
  calc_single_pos_gates = FALSE,
  debug = FALSE
)
\end{verbatim}
\end{Usage}
%
\begin{Arguments}
\begin{ldescription}
\item[\code{path\_project}] character. Path to project directory where all results will be saved.
This directory will contain subdirectories for each marker with gate tables,
statistics, and plots. The directory will be created if it doesn't exist.

\item[\code{.data}] GatingSet. A flowWorkspace GatingSet object containing the flow cytometry
data with both stimulated and unstimulated samples. The GatingSet should have
consistent channel names across all samples and include proper sample annotations.

\item[\code{batch\_list}] list. Named list where each element contains indices of samples
belonging to the same batch/donor. Names will be used for batch identification.
Example: list(donor1 = 1:10, donor2 = 11:20). Proper batching is crucial for
accurate background subtraction and gate identification.

\item[\code{chnl}] list. List where each element specifies parameters for gating a
specific channel Each element should be a list containing at minimum the channel
name (e.g., list(cut = "IL2")). Additional channel-specific parameters can
override global defaults. The marker name should match channel names in the GatingSet.

\item[\code{marker}] character vector. Alternative way to specify markers to gate on.
When provided, this is used instead of chnl to determine which markers to analyze.
Default is NULL.

\item[\code{pop\_gate}] character vector. Population(s) within which to perform gating.
Default is "root" to gate on all cells. Can specify other populations like
"CD3+" or "CD4+" if these gates already exist in the GatingSet.

\item[\code{bias\_uns}] numeric. Bias adjustment for unstimulated samples to account for
background cytokine production. When NULL (default), no bias correction is applied.
Positive values shift the unstimulated distribution higher, making gates more
conservative. Typically ranges from 0.1 to 1.0 when used.

\item[\code{bias\_uns\_factor}] numeric. Multiplicative factor applied to bias\_uns.
Default is 1. Values > 1 increase the bias effect, values < 1 decrease it.
This provides fine-tuning of the bias correction.

\item[\code{exc\_min}] logical. Whether to exclude minimum expression values during
analysis. Default is TRUE. Minimum values often represent technical artifacts
or compensation spillover and should typically be excluded.

\item[\code{cp\_min}] numeric. Minimum allowable cutpoint value. When NULL (default),
no minimum is enforced. Useful for ensuring gates don't fall below known
technical thresholds or background levels.

\item[\code{bw\_min}] numeric. Minimum bandwidth for density estimation. When NULL (default),
bandwidth is estimated automatically. Smaller values create more detailed density
estimates but may be noisier. Typical range is 0.01 to 0.1 on log-transformed data.

\item[\code{min\_cell}] numeric. Minimum number of cells required for reliable gating.
Default is 100. Samples with fewer cells will be skipped as they don't provide
sufficient statistical power for accurate gate identification.

\item[\code{max\_pos\_prob\_x}] numeric. Maximum x-value (expression level) to consider
when calculating positive probabilities. Default is Inf (no limit). Can be used
to exclude extremely high expression values that may represent doublets or artifacts.

\item[\code{gate\_quant}] numeric vector. Quantiles used for gate combination when multiple
gates are identified. Default is c(0.25, 0.75). The method specified in gate\_combn
determines how these quantiles are used (e.g., minimum of 25th percentiles).

\item[\code{tol\_clust}] numeric. Convergence tolerance for clustering algorithms used in
gate refinement. Default is 1e-7. Smaller values require more precise convergence
but may increase computation time.

\item[\code{gate\_combn}] character. Method for combining multiple gate candidates.
Default is "min" to use the most conservative (lowest) gate. Other options may
include "median" or "max" depending on the desired stringency.

\item[\code{marker\_settings}] list. Optional list of additional marker-specific settings
that override global defaults. Each element should be named with the marker name
and contain parameter overrides. Default is NULL.

\item[\code{chnl\_settings}] list. Optional list of channel-specific settings that override
global defaults. Similar to marker\_settings but keyed by channel names.
Default is NULL.

\item[\code{calc\_cyt\_pos\_gates}] logical. Whether to calculate refined cytokine-positive
gates using more sophisticated algorithms. Default is TRUE. When FALSE, only
basic gates are calculated, which may be less accurate but faster.

\item[\code{calc\_single\_pos\_gates}] logical. Whether to calculate single-positive gates
for individual markers in addition to combination gates. Default is FALSE.
Useful for detailed analysis of individual marker responses.

\item[\code{debug}] logical. Whether to enable detailed debug output and save intermediate
results. Default is FALSE. When TRUE, additional files and verbose output are
generated, useful for troubleshooting and method development.
\end{ldescription}
\end{Arguments}
%
\begin{Details}
The function implements a multi-step workflow for identifying cytokine-positive cells:

\strong{Step 1: Data Preparation}
\begin{itemize}

\item{} Validates input parameters and GatingSet structure
\item{} Completes marker specifications with default values
\item{} Organizes samples by batch for proper background subtraction

\end{itemize}


\strong{Step 2: Initial Gate Identification}
\begin{itemize}

\item{} Extracts expression data for each marker within specified populations
\item{} Estimates probability densities for stimulated and unstimulated samples
\item{} Identifies candidate cutpoints using outlier detection algorithms
\item{} Applies clustering to refine gate positions across batches

\end{itemize}


\strong{Step 3: Cytokine-Positive Gate Refinement (if calc\_cyt\_pos\_gates = TRUE)}
\begin{itemize}

\item{} Applies more sophisticated algorithms to refine initial gates
\item{} Accounts for background cytokine production and technical variability
\item{} Optimizes gates to minimize false positives while maintaining sensitivity

\end{itemize}


\strong{Step 4: Single-Positive Gates (if calc\_single\_pos\_gates = TRUE)}
\begin{itemize}

\item{} Calculates gates for individual markers independent of other markers
\item{} Useful for understanding single-marker responses

\end{itemize}


\strong{Step 5: Statistics Generation}
\begin{itemize}

\item{} Calculates comprehensive statistics including frequencies and combinations
\item{} Generates cross-tabulations of cytokine-positive populations
\item{} Saves results in structured format for downstream analysis

\end{itemize}


\strong{Important Considerations:}
\begin{itemize}

\item{} Ensure stimulated and unstimulated samples are properly paired by batch
\item{} Channel names in marker specifications must match GatingSet channels exactly
\item{} Sufficient cell numbers (min\_cell) are crucial for reliable gate identification
\item{} Background bias correction (bias\_uns) should be used cautiously and validated
\item{} Debug mode generates extensive output useful for method validation

\end{itemize}

\end{Details}
%
\begin{Value}
character. Returns the path to the project directory where all results
have been saved. The directory structure created includes:
\begin{itemize}

\item{} \code{path\_project/[marker\_name]/}: Directory for each marker containing:
\item{} \code{gate\_tbl\_init.rds}: Initial gate table with preliminary gates
\item{} \code{gate\_tbl.rds}: Final refined gate table
\item{} \code{stats/}: Directory containing statistics files
\item{} \code{plots/}: Directory containing visualization plots (if generated)

\end{itemize}

\end{Value}
%
\begin{SeeAlso}
\code{\LinkA{stimgate\_gate\_get}{stimgate.Rul.gate.Rul.get}} for extracting gate information,
\code{\LinkA{get\_stats}{get.Rul.stats}} for generating statistics from results,
\code{\LinkA{stimgate\_plot}{stimgate.Rul.plot}} for visualizing identified gates,
\code{\LinkA{stimgate\_fcs\_write}{stimgate.Rul.fcs.Rul.write}} for exporting cytokine-positive cells,
\code{\LinkA{GatingSet}{GatingSet}} for GatingSet documentation
\end{SeeAlso}
%
\begin{Examples}
\begin{ExampleCode}
{
example_data <- get_example_data()
gs <- flowWorkspace::load_gs(example_data$path_gs)
path_project <- file.path(tempdir(), "demonstration")

# Run gating
stimgate::stimgate_gate(
  .data = gs,
  path_project = path_project,
  pop_gate = "root",
  batch_list = example_data$batch_list,
  marker = example_data$marker
)

# Create plots
plots <- stimgate_plot(
  ind = example_data$batch_list[[1]], # indices in `gs` to plot
  .data = gs, # GatingSet
  path_project = path_project,
  marker = example_data$marker,
  grid = TRUE
)

# Advanced usage with parameter customization

}
\end{ExampleCode}
\end{Examples}
\HeaderA{stimgate\_gate\_get}{Get gates}{stimgate.Rul.gate.Rul.get}
%
\begin{Description}
Get all the gates for each of the markers gated.
\end{Description}
%
\begin{Usage}
\begin{verbatim}
stimgate_gate_get(path_project, pop = NULL, marker = NULL, chnl = NULL)
\end{verbatim}
\end{Usage}
%
\begin{Arguments}
\begin{ldescription}
\item[\code{path\_project}] character. Path to the project directory.

\item[\code{pop}] character. Optional population name(s) to filter gates by. Default is NULL (all populations).

\item[\code{marker}] character. Optional marker name(s) to filter gates by. Default is NULL (all markers).

\item[\code{chnl}] character. Optional channel name(s) to filter gates by. Default is NULL (all channels).
\end{ldescription}
\end{Arguments}
%
\begin{Value}
Gate table with gates for each sample for each marker.
\end{Value}
%
\begin{Examples}
\begin{ExampleCode}
{
# Get example dataset
example_data <- get_example_data()
gs <- flowWorkspace::load_gs(example_data$path_gs)

# Run the stimgate pipeline
path_project <- stimgate_gate(
  path_project = file.path(tempdir(), "get_gate_example"),
  .data = gs,
  batch_list = example_data$batch_list,
  marker = example_data$marker,
  pop_gate = "root"
)

# Get statistics for the identified gates
gates <- stimgate_gate_get(path_project)
}
\end{ExampleCode}
\end{Examples}
\HeaderA{stimgate\_meta\_read\_batch\_list}{Read batch list from project}{stimgate.Rul.meta.Rul.read.Rul.batch.Rul.list}
%
\begin{Description}
Read the saved batch\_list object from the project's meta\_data folder.
\end{Description}
%
\begin{Usage}
\begin{verbatim}
stimgate_meta_read_batch_list(path_project)
\end{verbatim}
\end{Usage}
%
\begin{Arguments}
\begin{ldescription}
\item[\code{path\_project}] character Path to project.
\end{ldescription}
\end{Arguments}
%
\begin{Value}
A list describing sample grouping into batches (as saved by .save\_meta\_data\_batch\_list()).
\end{Value}
%
\begin{Examples}
\begin{ExampleCode}
## Not run: 
tmp <- tempdir()
dir.create(file.path(tmp, "meta_data"), showWarnings = FALSE)
saveRDS(list(batch1 = c(1,2)), file.path(tmp, "meta_data", "batch_list.rds"))
stimgate_meta_read_batch_list(tmp)

## End(Not run)
\end{ExampleCode}
\end{Examples}
\HeaderA{stimgate\_meta\_read\_chnl\_lab}{Read channel or marker label mapping}{stimgate.Rul.meta.Rul.read.Rul.chnl.Rul.lab}
\aliasA{stimgate\_meta\_read\_marker\_lab}{stimgate\_meta\_read\_chnl\_lab}{stimgate.Rul.meta.Rul.read.Rul.marker.Rul.lab}
%
\begin{Description}
Read the saved channel label mapping (chnl\_lab.rds) from the project's meta\_data folder.
\end{Description}
%
\begin{Usage}
\begin{verbatim}
stimgate_meta_read_chnl_lab(path_project)

stimgate_meta_read_marker_lab(path_project)
\end{verbatim}
\end{Usage}
%
\begin{Arguments}
\begin{ldescription}
\item[\code{path\_project}] character Path to project.
\end{ldescription}
\end{Arguments}
%
\begin{Value}
Named character vector mapping channel names to labels.
\end{Value}
%
\begin{Examples}
\begin{ExampleCode}
## Not run: 
tmp <- tempdir()
dir.create(file.path(tmp, "meta_data"), showWarnings = FALSE)
saveRDS(c(BC1 = "BC1 label"), file.path(tmp, "meta_data", "chnl_lab.rds"))
stimgate_meta_read_chnl_lab(tmp)

## End(Not run)
\end{ExampleCode}
\end{Examples}
\HeaderA{stimgate\_meta\_read\_settings\_chnl}{Get marker settings for a single channel}{stimgate.Rul.meta.Rul.read.Rul.settings.Rul.chnl}
%
\begin{Description}
Retrieve the marker settings for a single channel. The function
accepts either a channel label (as returned by stimgate\_meta\_read\_chnl\_lab)
or the original channel name/key used in marker\_list.
\end{Description}
%
\begin{Usage}
\begin{verbatim}
stimgate_meta_read_settings_chnl(path_project, chnl)
\end{verbatim}
\end{Usage}
%
\begin{Arguments}
\begin{ldescription}
\item[\code{path\_project}] character Path to project.

\item[\code{chnl}] character Channel label or channel name.
\end{ldescription}
\end{Arguments}
%
\begin{Value}
A list of settings for the requested channel.
\end{Value}
%
\begin{Examples}
\begin{ExampleCode}
## Not run: 
tmp <- tempdir()
dir.create(file.path(tmp, "meta_data"), showWarnings = FALSE)
saveRDS(list(BC1 = list(a = 1)), file.path(tmp, "meta_data", "marker_list.rds"))
saveRDS(c(BC1 = "BC1 label"), file.path(tmp, "meta_data", "chnl_lab.rds"))
stimgate_meta_read_settings_chnl(tmp, "BC1 label")
stimgate_meta_read_settings_chnl(tmp, "BC1")

## End(Not run)
\end{ExampleCode}
\end{Examples}
\HeaderA{stimgate\_meta\_read\_settings\_chnls}{Read marker settings from project}{stimgate.Rul.meta.Rul.read.Rul.settings.Rul.chnls}
%
\begin{Description}
Read the saved marker settings list from the project's meta\_data folder.
\end{Description}
%
\begin{Usage}
\begin{verbatim}
stimgate_meta_read_settings_chnls(path_project)
\end{verbatim}
\end{Usage}
%
\begin{Arguments}
\begin{ldescription}
\item[\code{path\_project}] character Path to project.
\end{ldescription}
\end{Arguments}
%
\begin{Value}
A named list of marker settings (as saved by .complete\_chnl\_list\_save()).
\end{Value}
%
\begin{Examples}
\begin{ExampleCode}
## Not run: 
tmp <- tempdir()
dir.create(file.path(tmp, "meta_data"), showWarnings = FALSE)
saveRDS(list(BC1 = list(a = 1)), file.path(tmp, "meta_data", "marker_list.rds"))
stimgate_meta_read_settings_markers(tmp)

## End(Not run)
\end{ExampleCode}
\end{Examples}
\HeaderA{stimgate\_meta\_read\_settings\_marker}{Get settings for a named marker}{stimgate.Rul.meta.Rul.read.Rul.settings.Rul.marker}
%
\begin{Description}
Retrieve the settings for a marker by its original name/key.
\end{Description}
%
\begin{Usage}
\begin{verbatim}
stimgate_meta_read_settings_marker(path_project, marker)
\end{verbatim}
\end{Usage}
%
\begin{Arguments}
\begin{ldescription}
\item[\code{path\_project}] character Path to project.

\item[\code{marker}] character Marker name/key as stored in marker\_list.
\end{ldescription}
\end{Arguments}
%
\begin{Value}
A list of settings for the requested marker.
\end{Value}
%
\begin{Examples}
\begin{ExampleCode}
## Not run: 
tmp <- tempdir()
dir.create(file.path(tmp, "meta_data"), showWarnings = FALSE)
saveRDS(list(BC1 = list(a = 1)), file.path(tmp, "meta_data", "marker_list.rds"))
stimgate_meta_read_settings_marker(tmp, "BC1")

## End(Not run)
\end{ExampleCode}
\end{Examples}
\HeaderA{stimgate\_meta\_read\_settings\_markers}{Read marker list with channel labels}{stimgate.Rul.meta.Rul.read.Rul.settings.Rul.markers}
%
\begin{Description}
Read the project's marker list and return it with element names
replaced by channel labels (from chnl\_lab).
\end{Description}
%
\begin{Usage}
\begin{verbatim}
stimgate_meta_read_settings_markers(path_project)
\end{verbatim}
\end{Usage}
%
\begin{Arguments}
\begin{ldescription}
\item[\code{path\_project}] character Path to project.
\end{ldescription}
\end{Arguments}
%
\begin{Value}
A named list of marker settings where names are channel labels.
\end{Value}
%
\begin{Examples}
\begin{ExampleCode}
## Not run: 
tmp <- tempdir()
dir.create(file.path(tmp, "meta_data"), showWarnings = FALSE)
saveRDS(list(BC1 = list(a = 1)), file.path(tmp, "meta_data", "marker_list.rds"))
saveRDS(c(BC1 = "BC1 label"), file.path(tmp, "meta_data", "chnl_lab.rds"))
stimgate_meta_read_settings_chnls(tmp)

## End(Not run)
\end{ExampleCode}
\end{Examples}
\HeaderA{stimgate\_plot}{Plot stimulation gate}{stimgate.Rul.plot}
%
\begin{Description}
Plot bivariate hex and univariate density plots for batches of samples, along
with their gates.
\end{Description}
%
\begin{Usage}
\begin{verbatim}
stimgate_plot(
  ind,
  .data,
  path_project,
  marker = NULL,
  chnl = NULL,
  pop = NULL,
  ind_lab = NULL,
  axis_lab = NULL,
  exc_min = TRUE,
  limits_expand = NULL,
  limits_equal = FALSE,
  grid = TRUE,
  grid_n_col = 2,
  show_gate = TRUE,
  min_cell = 10
)
\end{verbatim}
\end{Usage}
%
\begin{Arguments}
\begin{ldescription}
\item[\code{ind}] numeric vector. Specifies indices in \code{.data} to plot.

\item[\code{.data}] GatingSet. Same GatingSet passed to \code{stimgate\_gate}.

\item[\code{path\_project}] character.
Path to the project directory used for \code{stimgate\_gate}.

\item[\code{marker}] character vector of length one or two. Specifies markers
to be plotted. If only one is passed, then only univariate plots are created.

\item[\code{chnl}] character vector of length one or two. Specifies channels
to be plotted. Ignored if \code{marker} is provided.

\item[\code{pop}] character. Specifies population within GatingSet that
gates were calculated on. If \code{NULL}, defaults to population specified
by folder name in \AsIs{\texttt{project\_path/gates/pop\_<pop>}}, but throws
an error if more than one population is detected (i.e. more
than one directory in \AsIs{\texttt{gates/}}). Default is \code{NULL}.

\item[\code{ind\_lab}] named character vector.
Labels for \code{ind} used in plot.
Optional.

\item[\code{axis\_lab}] named character vector.
Labels for axis titles, applied to \code{marker} or \code{chnl}.
Optional.

\item[\code{exc\_min}] Logical.
If \code{TRUE}, excludes the minimum expression values when processing the data.
Default is \code{TRUE}.

\item[\code{limits\_expand}] list.
Expand the limits of the plot axes.
Default is \code{NULL}.

\item[\code{limits\_equal}] Logical.
If TRUE, forces equal lengths of the limits.

\item[\code{grid}] Logical.
If TRUE, arranges the resulting plots in a grid format
using \code{cowplot::plot\_grid}.
Default is \code{TRUE}.

\item[\code{grid\_n\_col}] Integer.
Number of columns in grid layout.

\item[\code{show\_gate}] Logical.
If \code{TRUE}, overlays gate lines on the plots.|>
Default is \code{TRUE}.

\item[\code{min\_cell}] integer.
Minimum number of cells to be plotted.
Will skip plots with fewer cells.
Default is 10.
\end{ldescription}
\end{Arguments}
%
\begin{Value}
A grid of plots if \code{grid} is TRUE, otherwise a list of ggplot objects.
\end{Value}
%
\begin{Examples}
\begin{ExampleCode}
# Create example data and run gating
example_data <- get_example_data()
gs <- flowWorkspace::load_gs(example_data$path_gs)
path_project <- file.path(dirname(example_data$path_gs), "stimgate")

# Run gating
stimgate::stimgate_gate(
  .data = gs,
  path_project = path_project,
  pop_gate = "root",
  batch_list = example_data$batch_list,
  marker = example_data$marker
)

# Create plots
plots <- stimgate_plot(
  ind = example_data$batch_list[[1]], # indices in `gs` to plot
  .data = gs, # GatingSet
  path_project = path_project,
  marker = example_data$marker,
  grid = TRUE
)
\end{ExampleCode}
\end{Examples}
\printindex{}
\end{document}
